\documentclass{article}
\usepackage{amsmath,amssymb,hyperref}
\begin{document}
\section*{A finite obstruction to arbitrary SPD butterfly preconditioning}

This note records an exact finite computation, not a Lean theorem. It rules out
an auxiliary routing statement in which arbitrary symmetric positive definite
blocks replace the Cauchy blocks. It does not rule out the Cauchy statement,
weighted permutation tangent decoders, or an improvement for the original
attention model.

Let $H=2^d$. Index switches by triples $(k,s,g)$ with
$0\leq k<d$, $0\leq s<2^k$, and $0\leq g<2^{d-k-1}$, in that lexicographic
order. Given a bit $\theta_e$ at every switch, start with the original input
wire labels $0,\ldots,H-1$. At switch $e=(k,s,g)$, let $u_e,v_e$ be the
original labels currently at positions $s+2^{k+1}g$ and
$s+2^{k+1}g+2^k$. The row $B_e$ is $e_{u_e}^{\mathsf T}-e_{v_e}^{\mathsf T}$.
Swap those positions if $\theta_e=1$. Thus $B=B_\theta$ depends on the entire
chosen switch pattern.

For each node $(k,g)$ choose an SPD matrix $K_{k,g}$ of size $2^k$. Let $K$ be
the block operator that mixes the $s$ queries separately within each node.
The proposed auxiliary statement was that, for every $\theta$ and every such
choice of SPD blocks, there exist residues $a\in\mathbb R^H$ and a free bias
$\beta_{k,g}$ at each node such that
\[
 (1-2\theta_e)\bigl((KB_\theta a)_e+\beta_{k,g}\bigr)>0
 \quad\text{for every }e=(k,s,g).
\]

\paragraph{Exact counterexample.}
The attached file \texttt{LargeDomainSPDPreconditioner.json} specifies $d=7$,
all $448$ switch bits, and $19$ nonidentity blocks. Each listed block is given
as an integer matrix $R$ and means $K_{k,g}=RR^{\mathsf T}+I$; all other
blocks are the identity. Hence every block is SPD without a numerical spectral
calculation. It also specifies $144$ strictly positive integer coefficients
$\lambda_e$, with all unlisted coefficients zero. Put
$\mu_e=(1-2\theta_e)\lambda_e$. Direct integer calculation verifies
\[
 \sum_{s=0}^{2^k-1}\mu_{k,s,g}=0
 \quad\text{for every node},
 \qquad B_\theta^{\mathsf T}K\mu=0.
\]
The positive weighted sum, with weights $\lambda_e$, of the requested strict
inequalities is consequently zero. The system is infeasible, including all
$127$ free node biases.

The verifier \texttt{LargeDomainSPDPreconditionerVerify.py} uses only Python
integer arithmetic. It reconstructs the wire routing and applies each block
as $R(R^{\mathsf T}\mu)+\mu$. Discovery used floating linear programming to
find a support, followed by exact rational row reduction; the final verifier
does not trust either the floating solver or the row-reduction package.

\paragraph{Why the earlier tests were inconclusive.}
For $H=16$ there are only $32$ switch queries and at most $31$ residue and
node-bias dimensions. At $H=32$ there are $80$ queries and at most $63$ such
dimensions. Many random sign patterns can therefore be feasible even for a
model that is not universally feasible. Small random successes did not prove
an SPD routing theorem. At $H=128$, the exact certificate above settles the
arbitrary-SPD version negatively. A Cauchy proof, if one exists, must use more
than positive definiteness, such as the special total positivity or the
relation between its evaluation points.
\end{document}
